\documentclass{amsart}
% For dealing with references we use the comment environment
\usepackage{verbatim}
\newenvironment{reference}{\comment}{\endcomment}
%\newenvironment{reference}{}{}
\newenvironment{slogan}{\comment}{\endcomment}
\newenvironment{history}{\comment}{\endcomment}

% For commutative diagrams we use Xy-pic
\usepackage[all]{xy}

% We use 2cell for 2-commutative diagrams.
\xyoption{2cell}
\UseAllTwocells

% We use multicol for the list of chapters between chapters
\usepackage{multicol}

% This is generally recommended for better output
\usepackage{lmodern}
\usepackage[T1]{fontenc}

% For cross-file-references

% Package for hypertext links:
\usepackage{hyperref}

% For any local file, say "hello.tex" you want to link to please

% Theorem environments.
%
\theoremstyle{plain}
\newtheorem{theorem}[subsection]{Theorem}
\newtheorem{proposition}[subsection]{Proposition}
\newtheorem{lemma}[subsection]{Lemma}

\theoremstyle{definition}
\newtheorem{definition}[subsection]{Definition}
\newtheorem{example}[subsection]{Example}
\newtheorem{exercise}[subsection]{Exercise}
\newtheorem{situation}[subsection]{Situation}

\theoremstyle{remark}
\newtheorem{remark}[subsection]{Remark}
\newtheorem{remarks}[subsection]{Remarks}

\numberwithin{equation}{subsection}

% Macros
%
\def\lim{\mathop{\mathrm{lim}}\nolimits}
\def\colim{\mathop{\mathrm{colim}}\nolimits}
\def\Spec{\mathop{\mathrm{Spec}}}
\def\Hom{\mathop{\mathrm{Hom}}\nolimits}
\def\Ext{\mathop{\mathrm{Ext}}\nolimits}
\def\SheafHom{\mathop{\mathcal{H}\!\mathit{om}}\nolimits}
\def\SheafExt{\mathop{\mathcal{E}\!\mathit{xt}}\nolimits}
\def\Sch{\mathit{Sch}}
\def\Mor{\mathop{\mathrm{Mor}}\nolimits}
\def\Ob{\mathop{\mathrm{Ob}}\nolimits}
\def\Sh{\mathop{\mathit{Sh}}\nolimits}
\def\NL{\mathop{N\!L}\nolimits}
\def\CH{\mathop{\mathrm{CH}}\nolimits}
\def\proetale{{pro\text{-}\acute{e}tale}}
\def\etale{{\acute{e}tale}}
\def\QCoh{\mathit{QCoh}}
\def\Ker{\mathop{\mathrm{Ker}}}
\def\Im{\mathop{\mathrm{Im}}}
\def\Coker{\mathop{\mathrm{Coker}}}
\def\Coim{\mathop{\mathrm{Coim}}}

% Boxtimes
%
\DeclareMathSymbol{\boxtimes}{\mathbin}{AMSa}{"02}

%
% Macros for moduli stacks/spaces
%
\def\QCohstack{\mathcal{QC}\!\mathit{oh}}
\def\Cohstack{\mathcal{C}\!\mathit{oh}}
\def\Spacesstack{\mathcal{S}\!\mathit{paces}}
\def\Quotfunctor{\mathrm{Quot}}
\def\Hilbfunctor{\mathrm{Hilb}}
\def\Curvesstack{\mathcal{C}\!\mathit{urves}}
\def\Polarizedstack{\mathcal{P}\!\mathit{olarized}}
\def\Complexesstack{\mathcal{C}\!\mathit{omplexes}}
% \Pic is the operator that assigns to X its picard group, usage \Pic(X)
% \Picardstack_{X/B} denotes the Picard stack of X over B
% \Picardfunctor_{X/B} denotes the Picard functor of X over B
\def\Pic{\mathop{\mathrm{Pic}}\nolimits}
\def\Picardstack{\mathcal{P}\!\mathit{ic}}
\def\Picardfunctor{\mathrm{Pic}}
\def\Deformationcategory{\mathcal{D}\!\mathit{ef}}

\setlength{\emergencystretch}{3em}
\begin{document}
\title[Stacks Project, Tag 02TX]{451 --- Projective bundle formula for Chow homology of finite locally free modules}
\maketitle
\noindent Copyright (C) 2005--2025 Johan de Jong. Stacks Project authors. Source: \href{https://stacks.math.columbia.edu/tag/02TX}{Tag 02TX}.

\noindent Editorial difficulty note (not author text). Difficulty H2: Injectivity requires successively extracting every coefficient by Chern-class cap and proper pushforward. Surjectivity must be built by reduction to irreducible support and induction on the base dimension, followed by an affine-stratum filtration of the trivialized projective bundle and explicit lifts of cycle coefficients..

\noindent Editorial provenance note (not author text). History: excerpt from revision \texttt{a0444\allowbreak 6e57e\allowbreak c1fbc\allowbreak 252a8\allowbreak 71afc\allowbreak ec775\allowbreak 2fb28\allowbreak 07b14}, chow.tex, lines 6931--7059. Reference presentation was adapted; original-author context and core support blocks were retained or added. The mathematical wording is unchanged. Licensed under GNU FDL 1.2 or later; no invariant sections or cover texts. See ../licenses/COPYING. Context blocks reproduce author definitions and supporting results; they are not additional counted problems.

\begin{situation}
\label{situation-setup}
Here $S$ is a locally Noetherian, and universally catenary scheme.
Moreover, we assume $S$ is endowed with a dimension function
$\delta : S \longrightarrow \mathbf{Z}$.
\end{situation}

\begin{lemma}[Projective space bundle formula]
\label{lemma-chow-ring-projective-bundle}
Let $(S, \delta)$ be as in Situation \ref{situation-setup}.
Let $X$ be locally of finite type over $S$.
Let $\mathcal{E}$ be a finite locally free $\mathcal{O}_X$-module
$\mathcal{E}$ of rank $r$. Let $(\pi : P \to X, \mathcal{O}_P(1))$
be the projective bundle associated to $\mathcal{E}$.
The map
$$
\bigoplus\nolimits_{i = 0}^{r - 1}
\CH_{k + i}(X)
\longrightarrow
\CH_{k + r - 1}(P),
$$
$$
(\alpha_0, \ldots, \alpha_{r-1})
\longmapsto
\pi^*\alpha_0 +
c_1(\mathcal{O}_P(1)) \cap \pi^*\alpha_1
+ \ldots +
c_1(\mathcal{O}_P(1))^{r - 1} \cap \pi^*\alpha_{r-1}
$$
is an isomorphism.
\end{lemma}

\begin{proof}
Fix $k \in \mathbf{Z}$. We first show the map is injective.
Suppose that $(\alpha_0, \ldots, \alpha_{r - 1})$ is an element
of the left hand side that maps to zero.
By Lemma \href{https://stacks.math.columbia.edu/tag/02TW}{lemma cap projective bundle (Tag 02TW)} we see that
$$
0 = \pi_*(\pi^*\alpha_0 +
c_1(\mathcal{O}_P(1)) \cap \pi^*\alpha_1
+ \ldots +
c_1(\mathcal{O}_P(1))^{r - 1} \cap \pi^*\alpha_{r-1})
= \alpha_{r - 1}
$$
Next, we see that
$$
0 = \pi_*(c_1(\mathcal{O}_P(1)) \cap (\pi^*\alpha_0 +
c_1(\mathcal{O}_P(1)) \cap \pi^*\alpha_1
+ \ldots +
c_1(\mathcal{O}_P(1))^{r - 2} \cap \pi^*\alpha_{r - 2}))
= \alpha_{r - 2}
$$
and so on. Hence the map is injective.

\medskip\noindent
It remains to show the map is surjective.
Let $X_i$, $i \in I$ be the irreducible components of $X$.
Then $P_i = \mathbf{P}(\mathcal{E}|_{X_i})$, $i \in I$
are the irreducible components of $P$. Consider the commutative
diagram
$$
\xymatrix{
\coprod P_i \ar[d]_{\coprod \pi_i} \ar[r]_p & P \ar[d]^\pi \\
\coprod X_i \ar[r]^q & X
}
$$
Observe that $p_*$ is surjective. If $\beta \in \CH_k(\coprod X_i)$
then $\pi^* q_* \beta = p_*(\coprod \pi_i)^* \beta$, see
Lemma \href{https://stacks.math.columbia.edu/tag/02RG}{lemma flat pullback proper pushforward (Tag 02RG)}. Similarly for
capping with $c_1(\mathcal{O}(1))$ by
Lemma \href{https://stacks.math.columbia.edu/tag/02SU}{lemma pushforward cap c1 (Tag 02SU)}.
Hence, if the map of the lemma is surjective for each
of the morphisms $\pi_i : P_i \to X_i$, then the map is
surjective for $\pi : P \to X$. Hence we may assume $X$ is irreducible.
Thus $\dim_\delta(X) < \infty$ and in particular we may use
induction on $\dim_\delta(X)$.

\medskip\noindent
The result is clear if $\dim_\delta(X) < k$.
Let $\alpha \in \CH_{k + r - 1}(P)$.
For any locally closed subscheme $T \subset X$ denote
$\gamma_T : \bigoplus \CH_{k + i}(T) \to \CH_{k + r - 1}(\pi^{-1}(T))$
the map
$$
\gamma_T(\alpha_0, \ldots, \alpha_{r - 1})
= \pi^*\alpha_0 + \ldots +
c_1(\mathcal{O}_{\pi^{-1}(T)}(1))^{r - 1} \cap \pi^*\alpha_{r - 1}.
$$
Suppose for some nonempty open $U \subset X$ we have
$\alpha|_{\pi^{-1}(U)} = \gamma_U(\alpha_0, \ldots, \alpha_{r - 1})$.
Then we may choose lifts $\alpha'_i \in \CH_{k + i}(X)$ and we
see that $\alpha - \gamma_X(\alpha'_0, \ldots, \alpha'_{r - 1})$
is by Lemma \href{https://stacks.math.columbia.edu/tag/02RX}{lemma restrict to open (Tag 02RX)}
rationally equivalent to a $k$-cycle on $P_Y = \mathbf{P}(\mathcal{E}|_Y)$
where $Y = X \setminus U$ as a reduced closed subscheme.
Note that $\dim_\delta(Y) < \dim_\delta(X)$.
By induction the result holds
for $P_Y \to Y$ and hence the result holds for $\alpha$.
Hence we may replace $X$ by any nonempty open of $X$.

\medskip\noindent
In particular we may assume that $\mathcal{E} \cong \mathcal{O}_X^{\oplus r}$.
In this case $\mathbf{P}(\mathcal{E}) = X \times \mathbf{P}^{r - 1}$.
Let us use the stratification
$$
\mathbf{P}^{r - 1} = \mathbf{A}^{r - 1}
\amalg \mathbf{A}^{r - 2}
\amalg \ldots
\amalg \mathbf{A}^0
$$
The closure of each stratum is a $\mathbf{P}^{r - 1 - i}$ which is a
representative of $c_1(\mathcal{O}(1))^i \cap [\mathbf{P}^{r - 1}]$.
Hence $P$ has a similar stratification
$$
P = U^{r - 1} \amalg U^{r - 2} \amalg \ldots \amalg U^0
$$
Let $P^i$ be the closure of $U^i$. Let $\pi^i : P^i \to X$
be the restriction of $\pi$ to $P^i$.
Let $\alpha \in \CH_{k + r - 1}(P)$. By
Lemma \href{https://stacks.math.columbia.edu/tag/02TT}{lemma pullback affine fibres surjective (Tag 02TT)}
we can write $\alpha|_{U^{r - 1}} = \pi^*\alpha_0|_{U^{r - 1}}$
for some $\alpha_0 \in \CH_k(X)$. Hence the difference
$\alpha - \pi^*\alpha_0$ is the image of some
$\alpha' \in \CH_{k + r - 1}(P^{r - 2})$.
By Lemma \href{https://stacks.math.columbia.edu/tag/02TT}{lemma pullback affine fibres surjective (Tag 02TT)}
again we can write
$\alpha'|_{U^{r - 2}} = (\pi^{r - 2})^*\alpha_1|_{U^{r - 2}}$
for some $\alpha_1 \in \CH_{k + 1}(X)$.
By Lemma \href{https://stacks.math.columbia.edu/tag/02TR}{lemma relative effective cartier (Tag 02TR)}
we see that the image of $(\pi^{r - 2})^*\alpha_1$
represents $c_1(\mathcal{O}_P(1)) \cap \pi^*\alpha_1$.
We also see that
$\alpha - \pi^*\alpha_0 - c_1(\mathcal{O}_P(1)) \cap \pi^*\alpha_1$
is the image of some $\alpha'' \in \CH_{k + r - 1}(P^{r - 3})$.
And so on.
\end{proof}
\end{document}
